\section{Fault-Aware Mapping}\label{sec:mapping}

Placement decides which tier each component uses, but not which physical row of that tier. With stuck
cells this matters: a stuck high-order cell shifts its weight by a large fixed amount, so it does far
more harm under an important component than under an unimportant one. This section uses a property of
the SVD layer to choose the row of every component around a chip's stuck cells.

\subsection{Free Order of Components}
For any permutation matrix $P$,
\begin{equation}
(UP)\big(P^{\top}\Sigma V^{\top}\big)=U\Sigma V^{\top},
\label{eq:perm}
\end{equation}
so the $k$ components of a layer can be stored in any order without changing its output. We call the
$j$-th position of a tier a \emph{slot}: row $j$ of the tier's first crossbar together with column $j$ of
its second crossbar. Moving a component to another slot moves its row $\sigma_r v_r^{\top}$ and its
column $u_r$ together. The order is fixed when the chip is programmed, so a chosen mapping needs no
permutation hardware and no change to the neighbouring layers.

\subsection{Known Stuck Cells}
Stuck cells can be located after fabrication by writing and reading back every cell. We assume this test
has been done, so each chip's stuck-cell map is known before its weights are programmed; in simulation it
is the chip's drawn fault map. Programming noise is random and not known in advance, so the mapping
targets stuck cells only.

\subsection{Cost of a Component in a Slot}
Placing component $r$ in slot $s$ costs the expected rise in loss caused by the stuck cells of that slot.
A cell of slice $t$ stuck at conductance $G^{*}$, under an entry whose digit should be $d$, changes the
entry by
\begin{equation}
\delta=\pm\,s\,L^{t}\,\frac{G^{*}-(G_{\min}+d\,\Delta)}{\Delta},
\label{eq:stuckerr}
\end{equation}
with $s$ the 8-bit scale of the crossbar row holding the entry and the sign set by the side of the
differential pair. Weighting each error by the Fisher information of its entry, as in \cref{eq:fisher},
gives
\begin{equation}
C_{rs}=\sum_{\text{stuck cells }c\text{ in slot }s}F_{e(r,c)}\,\delta_{r,c}^{2},
\label{eq:cost}
\end{equation}
where $e(r,c)$ is the entry of component $r$ that would occupy cell $c$. The first-crossbar row moves with
its component and keeps its own scale $s_r$, while the scales $t_i$ of the second crossbar depend only on
which components share the tier, so every term of $C_{rs}$ is known before programming. The sum treats two
stuck cells under the same entry as independent, which slightly misestimates their joint error; it is
exact when each entry has at most one stuck cell.

\subsection{Assignment}
Within each tier, the mapping is the permutation $\pi$ that minimises $\sum_r C_{r\pi(r)}$, a linear
assignment problem. We solve it exactly with the Hungarian method~\cite{kuhn1955hungarian} as implemented
in SciPy~\cite{virtanen2020scipy}, separately for every layer and tier; cost matrices have at most
$614\times614$ entries (all components of a feed-forward layer). We apply it to the
magnitude hybrid, keeping its SLC set and choosing the slot of each component within its tier; this is
our \emph{fault-aware mapping}. As a high-cost reference, we apply the same mapping to an all-SLC chip.

\subsection{Joint Tier and Slot Mapping}
Fault-aware mapping keeps the SLC set fixed. A natural extension chooses tier and slot together for each
chip: an important component might be better off in a clean MLC slot than in a badly damaged SLC slot. We
form one assignment over all $k$ components and all $k$ slots of a layer, with $\lceil fk\rceil$ SLC slots,
which enforces the SLC budget. The cost of a slot adds the expected damage from programming noise in its
tier to \cref{eq:cost},
\begin{equation}
C^{\text{joint}}_{rs}=C_{rs}+\kappa_{\tau(s)}\Big(s_r^2\sum_j F\big(\sigma_r v_{rj}\big)+\sum_i t_i^2F\big(u_{ir}\big)\Big),
\label{eq:c1}
\end{equation}
where $\kappa_\tau=2(\sigma_d/\Delta_\tau)^2\sum_t L_\tau^{2t}$ is the noise variance of a weight in tier
$\tau$, in units of its 8-bit step. The second term is the Fisher score $D_r$ of \cref{eq:fisher} scaled
by the tier. We use $\sigma_d=\max(\sigma,2\%)$, so that on a chip with stuck cells but no programming noise
the tier choice still reflects the higher fragility of MLC.

\subsection{Checking the Cost Model}
Before every run, a self-test programs a random layer with a known fault map and compares the predicted
cost $\sum_r C_{rr}$ with the actual Fisher-weighted error of the programmed crossbar. The run stops unless
the two agree to within 10\% for both SLC and MLC, and unless the assignment does not increase the cost.
